%% toc: P2 9 9.5,9.6,9.7
\section{Proofs for the rules of \texorpdfstring{\(\mathsf S\)}{S} and the classical syllogisms}
\label{sec:annex-p02-w02c-p02-002-rules}

Derivations are those of Definition~\ref{def:p02-syllogistic-derivation}:
finite trees whose leaves are premises or instances of rules with no
antecedents, and whose other nodes are conclusions of rule instances whose
antecedents stand, up to the formula identifications of
Definition~\ref{def:p02-syllogistic-identifications}, at the immediate
predecessors.

\begin{annexitem}{annex:p02-w02c-p02-002-t-needed}{thm:p02-w02c-p02-002-t-needed}{Without T nothing is derivable from no premises}
\begin{proof}
\emph{Item 1.} The rules of \(\mathsf S\) and their numbers of antecedents
are: D1, D2, D3 and B with two each; A and I with one each; X with two; and T
with none. The rules of \(\mathsf S^\dagger\) are: D and B with two
antecedents each; A, I and N with one each; X with two; and T with none.
Thus in each sheet T is the only rule without antecedents, and every rule of
\(\mathsf S^-\) and of \(\mathsf S^{\dagger-}\) has at least one antecedent.

Suppose, for a contradiction, that some formula has a derivation from
\(\varnothing\) using only rules of \(\mathsf S^-\), or only rules of
\(\mathsf S^{\dagger-}\). The derivation is a finite tree, so it has at least
one leaf. By the definition of a derivation, that leaf is either a premise,
which is impossible because the premise set is empty, or an instance of a
rule with no antecedents, which is impossible because the rule set contains
no such rule. The formula identifications only allow an antecedent of a rule
instance to be matched with an identified formula at a predecessor; they do
not turn any formula into a leaf. Hence no such derivation exists.

\emph{Item 2.} Let \(\mathfrak A\) be any structure. Since
\(p^{\mathfrak A}\subseteq p^{\mathfrak A}\),
Definition~\ref{def:p02-syllogistic-semantics} gives
\(\mathfrak A\models\forall(p,p)\). As \(\varnothing\) is satisfied by every
structure, \(\varnothing\models\forall(p,p)\). The formula \(\forall(p,p)\) has
the \(\mathcal S\)-form \(\forall(p,l)\) with \(l=p\), and the
\(\mathcal S^\dagger\)-form \(\forall(l,m)\) with \(l=m=p\)
(Definition~\ref{def:p02-syllogistic-language}). By item~1 it has no
derivation from \(\varnothing\) in \(\mathsf S^-\) or in
\(\mathsf S^{\dagger-}\). So in each case a semantic consequence between
formulas of the fragment is not derivable, and the rule set is not complete
for the fragment.

\emph{Item 3.} In the sheet of \(\mathsf S\), T has no antecedents and
conclusion \(\forall(p,p)\) for an atom \(p\); in the sheet of
\(\mathsf S^\dagger\), T has conclusion \(\forall(l,l)\) for a literal \(l\),
and \(l=p\) is allowed. By the definition of a derivation, a T instance is a
permitted leaf. The tree consisting of that single leaf is therefore a
derivation of \(\forall(p,p)\) from \(\varnothing\) in each calculus.
\end{proof}
\end{annexitem}

\begin{annexitem}{annex:p02-w02c-p02-002-existential-closure}{thm:p02-w02c-p02-002-existential-closure}{Closure with an existential premise}
\begin{proof}
Throughout, \(\Phi=\{\exists(p,q)\}\). Call a universal formula
\emph{diagonal} if it is \(\forall(k,k)\) for some literal \(k\).
By Definition~\ref{def:p02-syllogistic-identifications}, \(\forall(k,k)\) is
identified only with \(\forall(\bar k,\bar k)\), which is again diagonal,
and an existential formula is identified only with an existential formula.

\emph{Item 1.} If \(l=m\), the T rule of \(\mathsf S^\dagger\) has the
instance \(\forall(l,l)\), and the one-node tree consisting of that instance
is a BTA derivation from \(\Phi\).

Conversely, take any derivation from \(\Phi\) using only B, T and A of
\(\mathsf S^\dagger\). We show by induction from the leaves that every node
carries either the existential premise \(\exists(p,q)\) or a diagonal
formula. A leaf is the premise \(\exists(p,q)\) or a T instance
\(\forall(k,k)\), which is diagonal. Consider a node that is the conclusion
of a rule instance, and assume the claim for its immediate predecessors.
\begin{itemize}
\item If the rule is T, the conclusion \(\forall(k,k)\) is diagonal.
\item If the rule is B, its antecedents \(\forall(l',m')\) and
 \(\forall(m',n')\) are universal. Each is identified with the formula at a
 predecessor, which is therefore universal, hence not \(\exists(p,q)\), hence
 diagonal by the induction hypothesis. A formula identified with a diagonal
 formula is diagonal, so \(l'=m'\) and \(m'=n'\). The conclusion
 \(\forall(l',n')\) is \(\forall(l',l')\), which is diagonal.
\item If the rule is A, its antecedent \(\forall(l',\bar l')\) is universal,
 so by the same reasoning it is diagonal, which requires \(l'=\bar l'\).
 This is impossible: if \(l'\) is an atom \(r\), then \(\bar l'\) is the
 negative literal \(\bar r\), and if \(l'=\bar r\), then \(\bar l'=r\); a
 positive literal is never a negative one. So no node is the conclusion of
 an A instance.
\end{itemize}
This proves the claim. If the derivation is a derivation of
\(\forall(l,m)\), its root carries a formula identified with
\(\forall(l,m)\). That formula is universal, so it is diagonal, and
therefore \(\forall(l,m)\) is diagonal: \(l=m\).

\emph{Item 2.} By item~1,
\begin{align*}
 V^\Phi&=\{m:\text{there is }l\in V\text{ with }
 \Phi\vdash_{\mathrm{BTA}}\forall(l,m)\}\\
 &=\{m:\text{there is }l\in V\text{ with }l=m\}=V.
\end{align*}
So \(V^\Phi\subseteq V\), that is, \(V\) is \(\Phi\)-closed. The set
\(\{\bar p\}\) is nonempty. It is consistent, since a pair \(l,\bar l\)
consists of two different literals and \(\{\bar p\}\) has only one member.

\emph{Item 3.} Let \(W_0\) be the set of all negative literals. It contains
\(\bar p\). It is consistent, because in every pair \(l,\bar l\) one literal
is positive and so lies outside \(W_0\). It is complete, because for every
literal \(l\) either \(l\) is negative and \(l\in W_0\), or \(l\) is an atom
\(r\) and \(\bar l=\bar r\in W_0\). By item~2 it is \(\Phi\)-closed. So a
consistent, complete, \(\Phi\)-closed extension of \(\{\bar p\}\) exists.

Now let \(W\) be any consistent, complete, \(\Phi\)-closed extension of
\(\{\bar p\}\). Define \(\mathfrak A\) on the universe \(\{a\}\) by
\(r^{\mathfrak A}=\{a\}\) if the atom \(r\) lies in \(W\), and
\(r^{\mathfrak A}=\varnothing\) otherwise. For an atom \(r\), consistency
and completeness put exactly one of \(r,\bar r\) in \(W\), so
\(a\in(\bar r)^{\mathfrak A}=\{a\}\setminus r^{\mathfrak A}\) if and only if
\(r\notin W\), if and only if \(\bar r\in W\). Hence the element \(a\)
satisfies exactly the literals in \(W\), and \(\mathfrak A\) is the
one-element structure that \(W\) defines. Because \(\bar p\in W\) and
\(W\) is consistent, \(p\notin W\), so \(p^{\mathfrak A}=\varnothing\).
Then \(p^{\mathfrak A}\cap q^{\mathfrak A}=\varnothing\), and
\(\mathfrak A\not\models\exists(p,q)\) by
Definition~\ref{def:p02-syllogistic-semantics}.
\end{proof}
\end{annexitem}
